% Supplemental software study. These are new floating-point models, not the
% frozen checkpoints used for integer certification or hardware measurements.
\subsection{Supplemental endpoint-pair generalization study}
\label{sec:pair-stage2}
To examine sensitivity to the partition and include another additive baseline,
we fitted new models on a fixed endpoint-pair split of the 211,043-row TON\_IoT
source file. Both traffic directions for an unordered pair of source and
destination IP addresses were assigned together; ports were excluded from the
group key. A fixed salted SHA-256 rule assigned pairs to training, validation,
and test buckets in proportions 60:20:20, yielding 119,334, 53,622, and 38,087
rows, respectively. No pair or complete 44-column row crosses these partitions.
This is not a host-, time-, or feature-disjoint split: training and test share
50 hosts, and 2,889 test rows (7.585\%) have transformed model inputs present in
training. The test contains 3,669 normal and 34,418 attack rows and lacks
backdoor, ransomware, and XSS examples.

All models share ten numeric features selected using binary training labels
only, and four categorical fields; preprocessing was fitted only on training. The configurations
were an additive degree-eight KAN, a depth-five decision tree, a 16-unit MLP,
and an additive cubic B-spline logistic GAM with five knots per numeric feature
and no interactions. Unseen categorical values use an all-zero one-hot block;
never-trained KAN category-zero entries were set to zero without changing
training logits. A single-seed validation pilot preceded the fixed five-seed
study; no model, hyperparameter, seed, or threshold was selected from its scores.
The five model seeds were 20260916--20260920; the preprocessing seed was fixed
at 20260916. All twenty fitted checkpoints were frozen before test evaluation,
and the decision threshold remained 0.5. These new floating-point checkpoints
are separate from the frozen integer exports and hardware cohorts.

\begin{table*}[t]
\centering
\caption{Supplemental endpoint-pair test: mean $\pm$ sample standard deviation
over five model seeds on the same 38,087 rows. Standard deviations describe
optimization variation on one fixed split; they are not confidence intervals.
FPR and TPR are percentages.}
\label{tab:pair-stage2}
\small
\begin{tabular}{lccccc}
\toprule
Model & BA & F1 & FPR (\%) & TPR (\%) & AUROC\\
\midrule
Additive KAN & $0.89484\pm0.00022$ & $0.98695\pm0.00009$ & $20.594\pm0.050$ & $99.563\pm0.019$ & $0.98488\pm0.00027$\\
DT5 & $0.68510\pm0$ & $0.96620\pm0$ & $62.687\pm0$ & $99.707\pm0$ & $0.49014\pm0$\\
MLP16 & $0.98885\pm0.00062$ & $0.99798\pm0.00015$ & $2.044\pm0.152$ & $99.815\pm0.039$ & $0.99806\pm0.00044$\\
Spline GAM & $0.97003\pm0$ & $0.99561\pm0$ & $5.724\pm0$ & $99.730\pm0$ & $0.99606\pm0$\\
\bottomrule
\end{tabular}
\end{table*}

The KAN exceeds DT5 in balanced accuracy on this split but remains below both
MLP16 and the spline GAM (Table~\ref{tab:pair-stage2}). Its F1 of 0.98695
coexists with a 20.594\% false-positive rate, which limits deployment usefulness.
For the 35,198 test rows whose transformed inputs are absent from training,
mean KAN BA is 0.89638, F1 is 0.98613, and FPR is 20.248\%; this subgroup has
3,623 normal and 31,575 attack rows. The matching-input subgroup contains only
46 normal rows, so its rates have limited normal-class support. These results
do not establish that additive structure improves detection or that the
canonical-split scores were inflated by leakage: the partition, class mix,
preprocessing fit, and model checkpoints differ. They support a narrower
contribution centered on verifiable integer deployment when an additive score
is required. The five seeds are not independent datasets, and neither the
validation pilot nor these test scores are pooled with the hardware cohort.
